% ============================================
% DOC-27 (EN)
% English translation — SPEC-935-R355
% ============================================
\section*{\textbf{DOC-27} --- The Origin (1853)}
\markboth{DOC-27 --- The Origin (1853)}{}

\noindent \textit{Document extracted from the expert file.}

\vspace{0.5em}\noindent\rule{\textwidth}{0.5pt}\vspace{0.5em}


\noindent {\footnotesize\texttt{PUBLIC ARCHIVE OF CEAR\'A $\cdot$ LOOSE DOCUMENT $\cdot$ 19TH CENT.}}\\[-0.2em]

\noindent\textbf{Document:} letter from Father José Inácio de Sousa to the President of the Province of Ceará, Public Relief Commission.\\
\noindent\textbf{Date:} March 14, 1853.\\
\noindent\textbf{Origin:} Parish of Nossa Senhora dos Remédios, Quixeramobim River.\\
\noindent\textbf{State of conservation:} paper stained by humidity; the lower edge bears a dark substance, unidentified, that soaked through every sheet of the bundle. Chemical analysis inconclusive.\\
\noindent\textbf{Attached by:} Dr. Lúcia Mendes, on February 12, 2026, with the note: \textquotedblleft{\textit{This proves the diary is not the beginning. The diary is only the echo.\textquotedblright{}}}

\noindent Most Excellent Sir President of the Province,

\noindent I take the liberty of writing to you because the living no longer hear me and the dead are too many to bury alone.

\noindent I have been here since November of last year, assigned to administer the sacraments to the refugees arriving on the riverbank. I received two hundred souls in January. Two hundred and forty in February. This month of March, I have lost count. They arrive in droves, entire families reduced to walking dust. Their eyes are the first thing that dies. Before the body. The eyes lose the light, the water, the life. They stay open, but they do not see. As if the drought had entered through them and dried the soul from within.

\noindent Those who die are taken to a trench half a league from the parish. I have no consecrated cemetery. I have no wood for coffins. I have no cloth for shrouds. I have only a long grave, opened in the dry earth, where the bodies are laid in layers. Earth over them. More bodies. More earth. The vultures have learned the way. They draw a circle in the sky that never breaks. Not even at night. They roost in the trees around it, waiting.

\noindent And it is about this that I must write to you.

\noindent For three nights now, I cannot say the Office of the Dead without feeling a presence. Not a spiritual presence, of the kind I am accustomed to discerning in thirty years of priesthood. It is a physical presence. It arrives with the smell.

\noindent The smell comes from the trench.

\noindent It is not the common smell of death, which I know well. It is a sweet smell. Sweeter than incense, sweeter than mandacaru blossom. A smell that enters through the nostrils and settles in the throat like a thick liquid. The refugees feel it too. They turn toward the trench and stand looking. They do not weep. They do not pray. They only look. As if they expected something to come out of there.

\noindent I myself went to the trench on the second night.

\noindent I carried a lantern, holy water, a crucifix. Priest's things. Things that have always worked. I reached the edge and lit the bottom.

\noindent What I saw were not bodies.

\noindent They were remains. Arms and legs rising from the earth like roots. Hands open upward, as if asking. A child's hand, small, with twisted fingers, pointed at the night sky as if accusing someone.

\noindent The smell was stronger there.

\noindent I knelt to bless the trench. I began the \textit{Requiem Aeternam}. On the third word, I felt the earth vibrate.

\noindent Not a tremor. A pulse. As if the ground had a heart.

\noindent I knelt. I pressed my right hand against the earth. It was warm. Not hot from the sun --- warm from within. The temperature of a sick man's body. The earth was feverish.

\noindent I buried my hand. Up to the wrist. The earth was soft, damp --- a thing impossible in this dry Sertão. I felt something brush my fingers. Something smooth. Soft. Something that moved at the touch, as if it were alive.

\noindent I pulled my hand out.

\noindent The earth closed by itself over the hole.

\noindent A man nearby --- a refugee named Vicente, who lost his wife and three children last week --- looked at me and said:

\noindent \textquotedblleft{\textit{Father, the trench is hungry.\textquotedblright{}}}

\noindent Your Excellency, I know what Your Lordship will think upon reading this. That isolation and hunger have shaken my judgment. That the heat and the dying have clouded my reason. I do not know if it is true. Perhaps it is. But I must record what I saw, so that, if I die here --- and I believe I will die here --- someone may know that this trench is no common trench.

\noindent On the third day --- yesterday --- I woke with the palm of my right hand marked.

\noindent A dark mark, rounded, in the center of the palm. It looks like a burn, but it does not hurt. The skin is intact. The mark simply... is there. As if someone had pressed a hot seal against my flesh while I slept.

\noindent I tried to wash it. It does not come off.

\noindent I tried to pray. The mark does not leave.

\noindent Perhaps it is only the earth that entered my skin. Perhaps it is the ulcer that is beginning. Perhaps it is the mind that is going.

\noindent But tonight, when I closed my eyes, I heard a voice.

\noindent Not a voice that comes from outside. A voice that comes from \textit{inside the earth}. It repeated the same word, without stopping, like one learning to speak:

\noindent \textquotedblleft{\textit{Hunger.\textquotedblright{}}}

\noindent \textquotedblleft{\textit{Hunger.\textquotedblright{}}}

\noindent \textquotedblleft{\textit{Hunger.\textquotedblright{}}}

\noindent The voice of a child. But no child survived to speak.

\noindent I ask leave to end this letter. The lantern is running low.

\noindent May God have mercy on us.

\noindent \textit{Father José Inácio de Sousa}\\
\noindent \textit{Parish of Nossa Senhora dos Remédios}\\
\noindent \textit{March 14, 1853}
\noindent \textit{Note from the Public Archive of Ceará: Father Sousa's letter was found in 1916, inside a parish registry book, during the renovation of the building. The priest died on March 15, 1853 --- one day after writing this letter. The cause of death was recorded as \textquotedblleft{}voluntary starvation\textquotedblright{}: the priest had stopped eating, giving his ration to the refugees. His body was buried in the same trench he describes.}
\noindent Attached to this document is the 1916 cataloging sheet, where an anonymous clerk wrote in pencil, in the margin: \textquotedblleft{\textit{The sweet smell never left this place. Transfer to the Public Archive of Minas Gerais. Case 1,260.\textquotedblright{}} The clerk was never identified. The note is in the same trembling handwriting that appears on the edges of Joaquim's diary.}




\noindent \textbf{July 2026 --- Belo Horizonte}\\
\noindent \textbf{Last professional record of Dr. Lúcia Mendes}

\noindent \textbf{Internal document --- Self-referral for clinical leave}\\
\noindent \textbf{Professional:} Lúcia Mendes, CRM-MG 28.391. \textbf{Date:} July 28, 2026.

\noindent \textbf{Ethical note:}

\noindent This record does not constitute a valid psychological report. It was written to justify immediate leave, interruption of research, and referral to independent assessment. The first person was kept because this is an urgent annotation.

\noindent \textbf{Summarized history:}
\noindent In January 2026, the professional came into contact with the diary attributed to patient 1,260, Joaquim Antônio Correia, in the context of research in historical forensic psychology. The material remained under intermittent custody, with subsequent failures of sealing and registration.

\noindent In the following weeks, persistent signs were observed after handling the archive:

\begin{itemize}
\item subjective change of coloration in the right eye, still without conclusive ophthalmological examination;
\item recurrent perception of a sweetish odor without an identified environmental source;
\item repetitive dreams of trench, ward, and blank form;
\item nighttime hypervigilance and difficulty concentrating;
\item impulse to return to the diary even after formal decision to end the analysis.
\end{itemize}

\noindent \textbf{Working hypotheses:}

\begin{itemize}
\item acute reaction to the traumatic material researched, with intrusive images and insomnia;
\item suggestion intensified by long documentary exposure;
\item countertransferential identification with patient 1,260;
\item mild dissociative phenomenon, to be confirmed by external assessment.
\end{itemize}

\noindent \textbf{Measures adopted:}

\begin{itemize}
\item suspension of interviews, academic writing, and circulation of copies;
\item referral of the original for judicial custody;
\item request for clinical assessment by a professional without prior contact with the case;
\item recommendation that the diary not be read individually.
\end{itemize}

\noindent \textbf{Unlogged annotation:}
\noindent The printed copy of this document bears, at the end, the number ``1,262'' written in pencil next to the professional's name. The stroke does not match the main signature. Below it there is an empty line, without identification.

\noindent \textit{Lúcia Mendes --- CRM-MG 28.391}\\
\noindent Record closed at 3:14 AM, according to the time noted at the bottom of the copy.




\vspace{0.5em}
\noindent\textit{\textbf{↪ Routes:} \rota{DOC-01} \rotasep \rota{LUC-02} \rotasep \rota{EPI-01}}